\scanpage{123}
à l'équateur, dans le sens de cette orientation.

On aura alors
\[
\deleted{T_{\widehat A_2}^{\unsure{Y}}=\mathbb U^A}
\]
\[
T_{J_{\mathrm I}}^Y=\mathbb U^{J_{\mathrm I}}\mathop{\wedge}_{\{\pm1\}}\omega_Q
\]
\margin{NB $\{\pm1\}$ s'envoie diagonalement dans $\mathbb U^{J_1}$}
\[
T_{J_{\mathrm {II}}}^Y\simeq\mathbb U^{J_2}
\]
\margin{car $\displaystyle\bigwedge_{\substack{i\in J_{\mathrm{II}}\\\text{sur }\alpha}}T_i\simeq\mathbb U$, puisque les $T_i$ sont isomorphes can. à $\mathbb U\mathop{\wedge}_{\pm1}\omega_Q$, donc le produit contracté de deux d'entre eux est $\mathbb U$}
\[
T_{J_{\mathrm{III}}}^Y\simeq
\prod_{i\in J_{\unsure{\mathrm{II}|\mathrm{III}}}}
\mathbb U\mathop{\wedge}_{\{\pm1\}}
\underbrace{(\omega_Q\wedge\widetilde J_{\mathrm{III},i})}_{\varepsilon_i}
\]

Donc, si $J_{\mathrm I}\ne\varnothing$, on introduit l'ensemble
\[
\unsure{\varepsilon|\mathcal E}_{G,Q}=\varepsilon_{G,Y_Q}
=\underbrace{\varepsilon_{\widehat A_1}}_{\displaystyle\prod_{a\in\widehat A_1}\varepsilon_a}
\times\omega_Q\times
\underbrace{\varepsilon_{J_{\mathrm{III}}}}_{\displaystyle\prod_{i\in J_{\mathrm{III}}}\omega_Q\wedge\widetilde J_{\mathrm{III},i}}
\]
qui est un torseur sous
\[
\gamma_{\unsure{G_s|G}}=\{\pm1\}^{\widehat A_1}\times\{\pm1\}\times\{\pm1\}^{J_{\mathrm{III}}}
\]
groupe qui s'envoie dans $\mathbb U^{\widehat A}$
\diagram{\begin{tikzcd}[column sep=large,row sep=small]
\gamma_G\arrow[r]\arrow[dr,bend right=18]&
\mathbb U^{\widehat A}\unsure{\simeq|=}\mathbb U^{\widehat A_1}\times
\mathbb U^{J_{\mathrm I}}\times
\mathbb U^{J_{\mathrm{II}}/\sigma}\times\mathbb U^{J_{\mathrm{III}}}\\
&\mathbb U^{\widehat A_1}\times\overset{\unsure{\smile|\vee}}{\mathbb U^{J_1}}\times\mathbb U^{J_{\mathrm{III}}}
\end{tikzcd}}
\deleted{\unsure{en f}}

en factorisant par le produit triple, et en

\mspage{122}
